\par\Needspace{6\baselineskip}\medskip
\begin{center}\hypertarget{native-245-article174}{}\pdfbookmark[2]{174.}{article-174}174.\end{center}
Ist die Bahn eines Himmelskörpers nach den in den vorangegangenen Abschnitten dargelegten Verfahren aus drei oder vier geozentrischen Positionen bestimmt worden, die ihrerseits nach der Vorschrift des vorigen Artikels jeweils aus mehreren Beobachtungen gewonnen worden waren. so wird diese Bahn gewissermaßen die Mitte zwischen all diesen Beobachtungen einnehmen, und in den Differenzen zwischen beobachteten und berechneten Orten wird keine Spur einer Regelmäßigkeit verbleiben, die sich durch Korrektur der Bahnelemente beseitigen oder merklich abschwächen ließe. Wenn sich die Gesamtheit der Beobachtungen über keinen allzu langen Zeitraum erstreckt, lässt sich auf diese Weise die gewünschte Übereinstimmung der Bahnelemente mit allen Beobachtungen erreichen, sofern man die drei oder vier sogenannten Normalpositionen geschickt wählt. Bei der Bestimmung der Bahnen von Kometen oder neuen Planeten, deren Beobachtungen sich noch nicht über ein Jahr hinaus erstrecken, kommen wir auf diese Weise gewöhnlich so weit, wie es die Natur der Sache erlaubt. Ist also die zu bestimmende Bahn in einem beträchtlichen Winkel gegen die Ekliptik geneigt, so wird man sie im Allgemeinen auf drei Beobachtungen stützen, die man so weit wie möglich auseinanderliegend wählt: Geraten wir dabei jedoch zufällig in einen der oben ausgeschlossenen Fälle (Art.\ 160---162) oder erscheint die Bahnneigung zu gering, so werden wir die Bestimmung aus vier Positionen vorziehen, die wir ebenfalls mit möglichst großem zeitlichen Abstand wählen.

Liegt dagegen bereits eine längere Beobachtungsreihe über mehrere Jahre vor, so lassen sich daraus mehr Normalpositionen ableiten: Der höchsten Präzision wäre also wenig gedient, wenn wir zur Bahnbestimmung nur drei oder vier Positionen auswählten und alle übrigen völlig außer Acht ließen. Vielmehr werden wir in einem solchen Fall, wenn wir die höchste Präzision anstreben, möglichst viele sorgfältig bestimmte Positionen zusammentragen und verwenden. Dann liegen also mehr Daten vor, als zur Bestimmung der Unbekannten erforderlich sind: Doch alle diese Daten sind mit Fehlern behaftet, mögen sie noch so gering sein, sodass es im Allgemeinen unmöglich ist, allen genau zu genügen. Nun gibt es keinen Grund, gerade sechs bestimmte dieser Daten als absolut genau zu betrachten. Nach den Grundsätzen der Wahrscheinlichkeit müssen wir vielmehr voraussetzen, dass in \SourceLeaf{246}allen Daten unterschiedslos größere oder kleinere Fehler gleichermaßen möglich sind; zudem treten im Allgemeinen geringere Fehler häufiger auf als größere: Daher entspricht eine Bahn, die sechs Daten genau genügt, von den übrigen aber mehr oder weniger abweicht, offenbar den Grundsätzen der Wahrscheinlichkeitsrechnung weniger als eine andere, die zwar auch von diesen sechs etwas abweicht, dafür aber mit den übrigen umso besser übereinstimmt. Die Untersuchung der Bahn mit der im strengen Sinne \emph{größten} Wahrscheinlichkeit hängt von der Kenntnis des Gesetzes ab, nach dem die Wahrscheinlichkeit der Fehler mit zunehmender Größe abnimmt: Diese Kenntnis hängt aber ihrerseits von so vielen vagen oder zweifelhaften Überlegungen ab --- darunter auch physiologischen ---, die sich keiner Rechnung unterwerfen lassen, dass man ein solches Gesetz in der praktischen Astronomie kaum je, wenn überhaupt, gehörig angeben kann. Gleichwohl ist die Untersuchung des Zusammenhangs zwischen diesem Gesetz und der wahrscheinlichsten Bahn, die wir nun in größter Allgemeinheit in Angriff nehmen, keineswegs als müßige Spekulation anzusehen.
